\documentclass[11pt]{article}
\usepackage[margin=1in]{geometry}
\usepackage{amsmath,amssymb,booktabs,hyperref}
\newcommand{\dd}{\mathrm d}
\title{A reciprocal carrier--medium density bridge}
\author{Zimmerman programme: dated candidate CPB-RB-2026-10-16}
\date{16 October 2026}
\begin{document}
\maketitle
\section{Candidate and scope}
This is a separately specified, nonrelativistic effective density-energy
candidate linking a cold carrier density to the existing nodal-fermion
response medium. It is not a change to the frozen theory, a derived
relativistic completion, or a matching calculation from the MeV scalar
action. Its new density coefficients are inputs. The result is constructive
at the local thermodynamic level: a reciprocal coupling with a sufficient
boundedness/convexity condition and a solved finite-charge benchmark.
The simplest strong-coupling branch also exposes a new obstruction to
preserving the desired galaxy law.

Density-mediated attraction and collapse in Bose--Fermi mixtures are
established subjects; Ref.~\cite{rb:mixture} is cited for that context,
not as a theorem about this candidate. The calculations below are derived
for the project's nodal energy and chosen contact terms. No global novelty
claim is made.

\section{Keep the occupation-dependent energy}
Let $\Delta=h|\mathbf P|\geq0$ be the nodal gap, and let $q$ denote
signed fermion number density relative to the specified filled negative
bands. This $q$ is not an electric charge density or the cold particle
count. In the zero-temperature low-energy patch, with the same quadratic
subtraction as earlier reports, the grand-potential density is
\begin{equation}
 \Omega_F(\Delta,\mu)=
 \begin{cases}
 A\Delta^3/3,& |\mu|\leq\Delta,\\
 A|\mu|\Delta^2/2-A|\mu|^3/6,& |\mu|>\Delta,
 \end{cases}\qquad A>0.
 \label{rb:grand}
\end{equation}
The $-A|\mu|^3/6$ term cannot be discarded when $\mu$ depends on
the carrier: it contributes to reciprocal density feedback. Differentiating,
\begin{equation}
 q=-\partial_\mu\Omega_F
 ={A\over2}\operatorname{sgn}\mu\,(\mu^2-\Delta^2)
                  \Theta(|\mu|-\Delta).
\end{equation}
Legendre transformation gives the particularly simple canonical energy
\begin{equation}
 F_F(\Delta,q)={A\over3}
                 \left(\Delta^2+{2|q|\over A}\right)^{3/2}.
 \label{rb:canonical}
\end{equation}
For $q\ne0$ define $E=\sqrt{\Delta^2+2|q|/A}$; then
\begin{equation}
 \partial_qF_F=\operatorname{sgn}q\,E,\qquad
 \partial_\Delta F_F=A\Delta E.
\end{equation}
At $q=0$ the charge subgradient is the interval $[-\Delta,\Delta]$.
The canonical formulation thus retains a neutral gapped region rather than
forcing a spurious smooth compressibility there.

These expressions require $\Delta,E,T\ll\Lambda$ and use the already
chosen critical subtraction. At finite cutoff there are analytic
corrections. They do not determine the filled band's absolute gravitational
energy or vacuum constant. A reference filled sea is additional matter
content; counting its signed excitations does not count its complete
cosmological inventory.

\section{An explicit reciprocal local interaction}
Let $n\geq0$ be the cold number density in this low-energy density model.
Choose
\begin{equation}
 \mathcal E_{\rm dens}=m_cn+{K_c\over2}n^2
       +\zeta nq+{u\over2}q^2+F_F(\Delta,q),
 \quad K_c>0,\quad u\geq{\zeta^2\over K_c}.
 \label{rb:energy}
\end{equation}
In natural units $[n]=[q]=E^3$, $[A]=E$, and
$[K_c]=[\zeta]=[u]=E^{-2}$. An operator-level motivation is a regulated
preferred-frame interaction
$\zeta\,\hat n_c:\!\Psi^\dagger\Psi\!:
+(u/2)(:\!\Psi^\dagger\Psi\!:)^2$, with carrier contact repulsion.
It preserves both number symmetries in this specified NR sector. The
density functional is a chosen Hartree-level description, not a rigorous
many-body bound for those regulated operators. Matching its $n,K_c,\zeta,u$
to the original $X,\Phi$ fields, including their neutral production stock
and number-changing channels, remains undone. In particular no condensate
or dilute coherent phase is assumed established for the MeV halo.

The sufficient stability condition follows from the exact identity
\begin{equation}
 {K_c\over2}n^2+\zeta nq+{u\over2}q^2
 ={K_c\over2}\left(n+{\zeta q\over K_c}\right)^2
  +{1\over2}\left(u-{\zeta^2\over K_c}\right)q^2.
 \label{rb:square}
\end{equation}
Also $F_F$ is jointly convex in $(\boldsymbol\Delta,q)$:
$|\boldsymbol\Delta|^2+2|q|/A$ is nonnegative convex and its $3/2$
power is increasing convex. On the smooth $q>0$ radial sector,
\begin{equation}
 F_{\Delta\Delta}=A(E+\Delta^2/E),\quad
 F_{qq}={1\over AE},\quad F_{\Delta q}={\Delta\over E},
 \qquad \det\operatorname{Hess}_{\Delta,q}F_F=1.
\end{equation}
The nonsmooth charge origin is covered by convexity/subgradients.
Consequently Eq.~\eqref{rb:energy} is nonnegative above the stated rest
term and has no negative local density-Hessian direction under the
sufficient condition. This is a local thermodynamic result for the
chosen patch energy, not full quantum, gradient, gravitational or dynamical
stability. With $u=0$ the coupled density Hessian can lose positivity;
a bare density coupling cannot simply inherit the old coldness proof.

Reciprocity is explicit. The chemical-potential shifts are
$\delta\mu_c=\zeta q$ and $\delta\mu_F=\zeta n+u q$.
For number currents $\mathbf j_c,\mathbf j_F$ with closed boundaries,
\begin{equation}
 {\dd\over\dd t}\int\zeta nq\,\dd^3x
 =\int\zeta(\mathbf j_c\cdot\nabla q+
                       \mathbf j_F\cdot\nabla n)\,\dd^3x.
\end{equation}
The carrier and fermion work terms are respectively
$-\zeta\mathbf j_c\cdot\nabla q$ and
$-\zeta\mathbf j_F\cdot\nabla n$, cancelling this change.
The $u q^2/2$ term has the corresponding internal pressure/work balance.
Thus the bridge transfers energy and momentum through an interaction;
it does not supply free cooling or transfer energy to a vacuum constant.
An actual occupation/transport evolution is still required.

\section{Reciprocal occupation and local force}
Use the existing positive-square field energy, once only,
\begin{equation}
 \mathcal E_{\rm field}={C\over2}
       |\nabla\varphi-\mathbf P/C|^2
       +F_F(h|\mathbf P|,q),\qquad C>0.
\end{equation}
Here $F_F$ is the same term as in Eq.~\eqref{rb:energy}; it is not added
twice and no extra bare cubic is included. A prescribed baryon source
has energy $q_b\rho_b\varphi$. Neglect gradients of $\mathbf P$ only
for this local-density static benchmark. In the spherical aligned branch,
\begin{equation}
 D=C\varphi'-P={q_bM_b(<r)\over4\pi r^2},\qquad
 {D\over C}=A h^2P E,
 \qquad \Delta E=S(r):={D\over CAh}.
 \label{rb:field}
\end{equation}
At fixed cold density, the fermion chemical equilibrium is
\begin{equation}
 \mu_F-\zeta n=\operatorname{sgn}q\,E+u q.
\end{equation}
When $|\mu_F-\zeta n|\leq\sqrt S$, the neutral solution has
$q=0$, $\Delta=E=\sqrt S$. Otherwise, with $s$ the sign of that
chemical potential, there is a unique $E>\sqrt S$ satisfying
\begin{equation}
 |\mu_F-\zeta n|=E+{uA\over2}(E^2-S^2/E^2),\quad
 \Delta={S\over E},\quad q={sA\over2}(E^2-\Delta^2).
\end{equation}
The right-hand side is strictly increasing. Charged regions satisfy
$\Delta\leq\sqrt S$: at fixed baryon flux the bridge suppresses this
local field response relative to the original neutral branch.

One first benchmark prescribed $n=1+2/(1+r^2)$ and solved
$\int q\,\dd^3x=0$ on $0.1\leq r\leq R$. Setting the background
chemical potential by hand to give only negative interior $q$ would violate
that inventory condition. The solution instead has inner negative charge,
a neutral cubic interval, and positive compensating outer charge.
Nine cases were checked by two Gauss grids and independently by adaptive
integration split at the occupation edges. This benchmark does not derive
the cold density. The next calculation removes that prescription.

\section{Vary both densities at fixed inventories}
Take a specified external baryonic potential energy per carrier $U(r)$.
Vary $n$, $q$ and the local scalar response with fixed cold number $N_c$
and fixed total fermion charge zero. On a positive-$n$ branch,
\begin{align}
 K_cn+\zeta q+U&=\nu,\qquad \nu=\mu_c-m_c,\\
 \operatorname{sgn}q\,E+u q+\zeta n&=\mu_F.
\end{align}
Eliminating $n$ gives
\begin{equation}
 \operatorname{sgn}q\,E+u_{\rm eff}q
       =\widetilde\mu+{\zeta U\over K_c},\quad
 u_{\rm eff}=u-{\zeta^2\over K_c}\geq0,\quad
 \widetilde\mu=\mu_F-{\zeta\nu\over K_c}.
 \label{rb:eliminate}
\end{equation}
This feedback is absent if the cold density is frozen. If the solution
crosses $n=0$, a constrained active-set problem replaces these formulas;
the executed examples explicitly check $n>0$.

Choose dimensionless $A=C=h=K_c=1$, $u=\zeta^2+1$,
\begin{equation}
 U(r)=-{2\over\sqrt{1+r^2}},\qquad
 D(r)={r\over(1+r^2)^{3/2}},\qquad 0.1\leq r\leq R.
\end{equation}
These are a prescribed external Plummer trap and scalar baryon flux, not
a self-consistent gravitational halo. The cold inventory is fixed at the
uncoupled value $N_c=\int(1-U)\,\dd^3x$. Because $\int q=0$,
the integrated first equation gives $\nu=1$ exactly and
\begin{equation}
 n(r)=1-U(r)-\zeta q(r),\qquad
 \int[n-(1-U)]\,\dd^3x=-\zeta\int q\,\dd^3x=0.
\end{equation}
The local cold profile is now an output of the reciprocal equilibrium.
For $\zeta=1,u=2$, adaptive results are:
\begin{center}
\begin{tabular}{rrrrr}
\toprule
$R$ & $\widetilde\mu$ & $Q_F^+=|Q_F^-|$ & neutral interval
 & maximum $\delta n/n_{\rm uncoupled}$\\
\midrule
10 & .344375 & 6.03679 & $[2.8165,8.6441]$ & .1913\\
30 & .122857 & 41.4106 & $[8.1088,24.3947]$ & .2307\\
100 & .0384464 & 208.539 & $[26.0006,78.0232]$ & .2461\\
\bottomrule
\end{tabular}
\end{center}
The full runs include $\zeta=0.3,1,3$. Both inventories and all three
stationarity equations are checked. Every sample has positive cold density.
There is also a simple uniform check for the selected roots:
$0<\widetilde\mu<1$, $0<\zeta\leq3$ and $u_{\rm eff}=1$.
In positive-charge regions $q+\sqrt{2q}\leq\widetilde\mu<1$, so
$q<(\sqrt3-1)^2/2$. Negative-charge regions increase $n$.
Consequently $n>1-3(\sqrt3-1)^2/2>0.196$ throughout the annulus,
not just at quadrature nodes. The numerical roots satisfy the assumed
chemical-potential interval with a finite margin.
Furthermore,
the sampled excitation energies remain below one percent of the chosen
cutoff $\Lambda=1000$. The cutoff test does not bound omitted nonlocal
gradient terms or finite-temperature effects. The container and each
inventory are finite for every example. Reflecting boundaries are external
constraints here; a physical wall/interface stress and energy are not
included. The cold inventory $N_c$ changes across the
different-$R$ family. None is a formed galaxy or a physical dimensional fit.

\section{A proved boundary-sensitive limit}
There is a discriminating asymptotic for this benchmark. Write the large-$r$
inputs as $U\sim-V_0/r$ and $S\sim D_0^2/r^2$, and define
$B=\zeta V_0/K_c>D_0$. Consider the family of larger annuli with the
same background density and fixed zero fermion charge. Then
\begin{equation}
 \widetilde\mu R\longrightarrow2B,\qquad
 {r_{\rm neutral,in}\over R}\longrightarrow{1\over2}-{D_0\over2B},
 \qquad
 {r_{\rm neutral,out}\over R}\longrightarrow{1\over2}+{D_0\over2B}.
 \label{rb:boundary}
\end{equation}
To prove this, set $r=Rx$, $\widetilde\mu=a/R$ and $w=ax-B$.
At fixed $x>0$, the $u_{\rm eff}q$ term is subleading. In charged regions,
\begin{equation}
 E\sim{|w|\over Rx},\quad
 \Delta\sim{D_0^2\over Rx|w|},\quad
 q\sim{A\,\operatorname{sgn}w\over2R^2x^2}
          (w^2-D_0^4/w^2),\qquad |w|>D_0.
\end{equation}
The neutral interval is $|w|\leq D_0$. After multiplying by the spherical
volume element, $Q_F/R$ tends to
\begin{equation}
 {2\pi A\over a}\{J(a-B)-J(B)\},\qquad
 J(t)={t^3\over3}+{D_0^4\over t}-{4D_0^3\over3},\quad t>D_0,
\end{equation}
in a neighborhood of its root $a=2B$.
The derivative $J'(t)=t^2-D_0^4/t^2$ is strictly positive.
Thus neutrality uniquely selects $a=2B$, giving Eq.~\eqref{rb:boundary}.
The limit is controlled: $E\leq|\widetilde\mu+\zeta U/K_c|$ in
charged regions and $|U(r)|\leq V_0/r$ give
$R^2x^2|q(Rx,a/R)|\leq A(|a|x+B)^2/2$ on the annulus.
Extend the integrand by zero below its inner endpoint. This is an
integrable bound on $0\leq x\leq1$, uniform in $R$ for fixed $a$.
Dominated convergence and brackets on either side of
the strictly monotone charge root give the stated convergence, rather
than presuming a limit of the roots.

For $D_0=1$, $\zeta=1$ has $B=2$. At $R=10000$ the calculation gives
$\widetilde\mu R=3.99499$ versus limit $4$, and neutral edge fractions
$(.250314,.750941)$ versus $(.25,.75)$. For $\zeta=3$, $B=6$,
the corresponding values are $11.9660$ versus $12$, and
$(.417849,.584989)$ versus $(5/12,7/12)$.
The $R=1000$ comparisons approach those limits less closely. These
finite checks support the implementation; the asymptotic follows from the
displayed scaling and charge integral.

This is an increasingly large background-bath family, not a way to give a
single isolated galaxy infinite cold stock. It proves that the neutral
cubic interval can track the boundary rather than converge to a fixed
interior region. For the ordered limit $R\to\infty$ followed by a large
fixed radius with $B>D_0$, the interior charged branch has
$E\sim B/r$ and $\Delta\sim D_0^2/(Br)$.
If both the prescribed gravitational trap strength $V_0$ and scalar source
amplitude $D_0^2$ are proportional to $M_b$ at fixed couplings, this
coefficient $D_0^2/B$ becomes independent of $M_b$. It does not give the
required $\sqrt{M_b}$ scaling. This is a conditional obstruction to the
strong minimal density bridge as a universal mass--velocity selector,
not a theorem excluding weaker couplings, different states, self-gravity,
nonlocal response or other bridge operators.

\section{Mediator matching: a constructive option and a charge gate}
An ordinary massive vector provides a useful microscopic test of the
coefficient matrix. Let it couple to conserved currents as
$V_\mu(g_cJ_c^\mu+g_FJ_F^\mu)$, with a healthy positive mass $m_V$.
The static Gauss operator is $-\nabla^2+m_V^2$ and its inverse is
positive. At wavelengths much larger than $m_V^{-1}$, the density energy is
\begin{equation}
 \delta\mathcal E_V={1\over2m_V^2}(g_cn_Q+g_Fq)^2.
\end{equation}
Here $n_Q=J_c^0$ is the signed cold charge density. With bare repulsions
$K_0,u_0\geq0$, the matched coefficients are
\begin{align}
 K_c&=K_0+g_c^2/m_V^2,&
 u&=u_0+g_F^2/m_V^2,&
 \zeta&=g_cg_F/m_V^2,\\
 K_cu-\zeta^2&=K_0u_0+
              (K_0g_F^2+u_0g_c^2)/m_V^2\geq0.
\end{align}
The chosen $K_c=1,u=\zeta^2+1$ matrix is realizable in this contact
sector with $u_0=0$, $K_0=1/(\zeta^2+1)$,
$g_c^2/m_V^2=\zeta^2/(\zeta^2+1)$ and
$g_F^2/m_V^2=\zeta^2+1$. Thus the positive coefficient relation need
not be an unexplained cancellation of unrelated negative terms.
This statement does not supply a causal relativistic completion of the
preferred-frame nodal bands or fix physical mediator masses/couplings.

The charge identification is decisive. Gauge invariance of
$X^*\Phi^2$ requires $Q_X=2Q_\Phi$. Its leading NR current is
\begin{equation}
 n_Q=n_{\Phi+}-n_{\Phi-}+2(n_{X+}-n_{X-}),
\end{equation}
whereas the positive cold stock includes the sum of particle and antiparticle
densities. The Higgs process $h\to XX^*\to2\Phi+2\Phi^*$ creates
zero total cold charge. On its equal-sign-density branch $n_Q=0$ even
with substantial positive rest mass. One therefore cannot replace $n_Q$
by the positive $n$ in the preceding benchmark and call the vector a
completed bridge to that neutral stock. A charged subcomponent/asymmetry
or a separately solved charge-polarization dynamics is required.
Neutrality does not forbid charge fluctuations or particle scattering,
but it removes this
identification of the mean Hartree source. No charged production history
has been supplied here.

A charge-even alternative is a real scalar mediator with an energy term
\begin{equation}
 {m_\sigma^2\over2}\sigma^2+\sigma(a_cn+a_Fq).
\end{equation}
Eliminating its static homogeneous value gives
$\delta\mathcal E_\sigma=-(a_cn+a_Fq)^2/(2m_\sigma^2)$.
The attractive sign differs from the vector's sign. With bare repulsions,
\begin{align}
 K_c&=K_0-a_c^2/m_\sigma^2,&
 u&=u_0-a_F^2/m_\sigma^2,&
 \zeta&=-a_ca_F/m_\sigma^2,\\
 K_cu-\zeta^2&=K_0u_0-
              (K_0a_F^2+u_0a_c^2)/m_\sigma^2.
\end{align}
Positive final diagonal coefficients and nonnegative determinant are
needed; scalar exchange alone does not generate the stabilization.
A coupling $\sigma|\Phi|^2$ is charge-even and respects the original
dark $U(1)$. Coupling the same scalar to signed medium density uses
$\sigma u_\mu J_F^\mu$, where $u^\mu$ is the medium rest-frame vector.
This is an allowed preferred-frame EFT possibility, with additional bare
contacts and a medium-state assumption, not an established covariant
particle theory. Its masses, repulsions, matching, finite-temperature
response and energy receiver remain to be calculated. Both mediator
options preserve the distinction between neutral cold mass and signed
fermion/cold charge rather than hiding it in a density symbol.

\section{What this supplies and the next obligation}
The previously unspecified bridge is now an explicit effective density
interaction with reciprocal chemical forces, a sufficient convexity
condition, a counted signed fermion inventory, and a jointly responding
cold population in a specified external potential. These are stronger
results than keeping a prescribed cold density or a fixed chemical-potential
reservoir. The added $u$ term and all bridge coefficients remain new inputs;
they are not explained by the original particle model or dark energy.

The next decisive calculation is to match a charge-even bridge to the
neutral particle action, or supply and count an actual charged subcomponent,
then determine its state/temperature and transport,
then solve the nonlocal occupied-band problem with live gravitational
sources and a physical boundary. A vacuum response term, an energy receiver
and absolute metric stress must be counted once. The present energy transfer
is internal to cold/fermion/field sectors and has no vacuum $Q$ by itself.
The source-assisted Higgs-production hierarchy remains independently open.
No galaxy fit, lensing prediction, origin of $\kappa$ or vacuum magnitude,
or whole-theory closure follows from these examples.

\section{Evidence}
The deterministic scripts and manifests are
\texttt{reciprocal\_density\_bridge\_run\_001} (151 controls),
\texttt{bridge\_charge\_resolution\_run\_001} (88),
\texttt{joint\_cold\_bridge\_run\_001} (97), and
\texttt{bridge\_large\_radius\_run\_002} (46), and
\texttt{bridge\_mediator\_matching\_run\_001} (16 symbolic identities).
Counts include inherited symbolic and local checks and are not mutually
independent physics validations. The first large-radius run completed the
mathematics but failed to collect the nested namespace's check list for
output; the corrected second run changes that bookkeeping only. No gate
was weakened. Proof and scope review are self-review, not an independent
referee or novelty certificate.

\begin{thebibliography}{9}
\bibitem{rb:mixture} K. Manabe and Y. Ohashi,
Thermodynamic stability, compressibility matrix, and effects of mediated
interactions in a strongly-interacting Bose--Fermi mixture,
Phys. Rev. A 103 (2021) 063317,
\href{https://arxiv.org/abs/2103.16843v1}{arXiv:2103.16843v1}.
The primary abstract authenticates the established density-collapse and
repulsive-stabilization context. No quantitative result for the present
nodal model is imported from it.
\end{thebibliography}
\end{document}
